\SourceLeaf{320}%
that is satisfied in the present investigation. It is only necessary to note that \(\psi^0\) and \(\psi'\) there are exactly \(T+\frac12t-(U+\frac12u)\) and \(T-\frac12t-(U-\frac12u)\) here, that \(h\) there is the same as \(s\) here, and finally that \(\chi\) there can in general differ from the quantity denoted here by \(U\) only by a quantity of second order, though in any case the difference is of at least first order. This gives, accurately through the third order inclusive,
\[
\begin{aligned}
T+\tfrac12t&=U+\tfrac12u-\frac{ee\cos P\mathord{.}\sin P\mathord{.}\sin U\mathord{.}\cos U^2}{12\cos\varphi\cos\theta}\mathord{.}s^3\\
T-\tfrac12t&=U-\tfrac12u+\frac{ee\cos P\mathord{.}\sin P\mathord{.}\sin U\mathord{.}\cos U^2}{12\cos\varphi\cos\theta}\mathord{.}s^3
\end{aligned}
\]
and therefore, with the same accuracy,
\[
\begin{aligned}
T&=U\quad .\ .\ .\ .\ .\ .\ .\ .\ .\ .\ .\ .\ .\ .\ .\ .\ .\ .\ .\ .\quad(6)\\
t&=u-\frac{ee\cos P\mathord{.}\sin P\mathord{.}\sin U\mathord{.}\cos U^2}{6\cos\varphi\cos\theta}\mathord{.}s^3\quad .\ .\ .\ .\ .\ .\ .\ .\ .\quad(7)
\end{aligned}
\]

\indent The comparison between the length of the geodesic line on the ellipsoid and the great-circle arc on the sphere was not worked out specifically in Art. 15 for the case under consideration, but it is very easy to supply it. In the notation used there, the length of the geodesic arc takes the form
\[
=A\int\frac{\cos y}{m\cos\psi}\mathord{.}\mathrm{d}x
\]
where the integration is to extend from \(x=-\frac12(h-\delta)\) to \(x=+\frac12(h+\delta)\). Since \(y\) and \(\psi\) are quantities of only third order, it is easy to see that omitting the factor \(\frac{\cos y}{\cos\psi}\) can introduce only an error of seventh order in the value of the integral. The length is therefore, accurately through the fifth order inclusive,
\[
\begin{aligned}
=A\int\frac{\mathrm{d}x}{m}&=A\int\mathrm{d}x(1-\mu x^3-\mu'x^4)\\
&=A(x-\tfrac14\mu x^4-\tfrac15\mu'x^5)+\text{Const.}\\
&=A\{h-\tfrac18(h^3\delta+h\delta^3)\mu-\tfrac1{80}(h^5+10h^3\delta\delta+5h\delta^4)\mu'\}
\end{aligned}
\]
The coefficients \(\mu,\mu'\) can be obtained by taking the series
\[
m=1-\frac{2ee\cos P\mathord{.}\sin P}{3\cos\varphi\cos\theta}\mathord{.}q^3-\frac{ee\cos P^2}{6\cos\varphi^2\cos\theta^2}(1-7ee\sin P^2)q^4\quad .\ .\ .\ .
\]
(which follows directly from the one given in Art. 9), making the substitution for \(q\)
